\section{Reference profiles and encodings}
\label{app:profiles}
The artifact names three existing profiles: D1 for dynamic delegation, S1 for
additive swarm evolution and F1 for funded work. The composed receiver joins
their original records; it does not replace them with another wire format.

\subsection{Delegation records}
A D1 allocation digest binds the persistent allocator namespace, immutable root
slot and target slot. Each holder-signed subdivision binds the parent's
allocation digest. The receiver's offer binds the allocation, complete slot,
receiver, effect, input, price ceiling and expiry. The holder's selection
includes the complete signed offer, capability digest, request ID and expected
selection revision. The allocator-signed permit carries the frozen slot and
selection. Signed domains are versioned; the qualified records use v2.

The original receiver envelope contains a slot ID, payload and permit.
Its input digest binds the slot ID and exact payload. In the composed layout,
the permit is carried in the governed request context and the payload is the
entire original tool-argument object. The receiver selects the installed layout;
the caller cannot downgrade it. The enclosing permit has its own signature and
is excluded from the input digest to avoid a hash cycle.

Canonical numbers are bounded by $2^{53}-1$, delegation depth by 32, children per
slot by 64, slots per root by 4,096 and envelopes by 64\,KiB. Acceptance is a
conjunction of 1--16 JSON-pointer equalities or inclusive integer ranges.
The one-invocation native capability is issued by the receiving kernel and
names the exact subject, server and tool. Per-call and total monetary ceilings
must be no greater than the signed offer, in the same currency.

The D1 offer specifies accepted-output pricing. The installed guard rechecks
the binding before dispatch and validates the exact JSON value after output
transformations. A qualifying predicate rejection withholds output and uses
the qualified reversible-hold path for a zero-charge denial. The executed
invocation remains consumed. Native capability validity and revocation remain
independent requirements.

Sealing freezes the selection, then persists its signed permit before returning
it. A crash between those durable steps can strand capacity. When a stored
permit is available, a retry returns its original issuance time and signature.
There is no post-seal replacement or automatic reclaim operation.

\subsection{Swarm and receiver records}
S1 retains the graph envelope, existing nodes, edges, joins, routes, witness
chains, allocations, pool and revocation epoch. Existing continuations retain
every field except the new graph digest and signature. New tokens and
allocations name new tasks; each task has at most one of each.

The SQLite installer verifies both bundles at store-owned time and commits the
head transition and parent archive together. Historical lookup checks the
graph ID, signed graph digest, full bundle digest and strict encoding.
Other stores need equivalent protected serialization and historical binding
to implement the same profile.

The composed receiver pins its allocator, witness, treaty peer, program and
runtime profile in protected deployment policy. Each open installs the runtime
hook and delegation guard before reconciliation. The existing runtime database
must remain present; losing it cannot provision an empty continuation domain.
The funded journal's policy hash prevents removal of the composed configuration.

Bilateral admission uses the existing signed treaty, permission antecedents,
invocation and continuation records. Such antecedents establish permission;
completed-work evidence comes from the retained native outcome. Execution export
revalidates that outcome and its frozen federation provenance under the
original participant identities and admission time. It grants no new dispatch,
settlement or remote-completion authority.

\subsection{Funding and deadlines}
F1's native agreement authenticates the full final request, original receiver
authority, checker context and settlement terms. The current workload uses
100 mock units per native agreement. The artifact-only F1 agreement is a
different concrete schema; the companion maps both explicitly.

Let $t_s<t_c<t_r<t_f$ be submission, challenge, resolution and refund boundaries.
Submission is allowed by $t_s$, first verifier decision recording when
$t_c<t\leq t_r$, and undecided refund when $t>t_f$. The interval before decision
is a timing boundary; the profile has no public challenge court.
An available failed check yields rejection; missing required output or custody
yields no eligible decision.

Acceptance is a domain-bound EIP-712 decision over the allocation, agreement,
submission, decision digest, acceptance bit, beneficiary, asset and amount.
The reference escrow has no later expiry for an earned withdrawal. Its
administrative pause and verifier-registry updates govern new funding, without
rewriting an existing claim. This fixed-price profile has no partial-settlement
entrypoint. Fees, disclosure loss and opportunity costs fall outside its
reserved-unit bound unless separately provisioned.
